\chapter{Score-Based Diffusion Models}
\label{app:score_diff_models}

\subsection{Score Matching}
\label{app:subsec:sm}

Chapter \ref{ch:literature_review}  introduced the Fisher divergence as the natural objective for
learning the score function, noting that it is intractable because it requires
$\nabla_x \log p(x)$. Hyv\"arinen \cite{hyvarinen2005estimation} showed that
the Fisher divergence can be rewritten in a form computable without the true
score. Starting from:
\[
    J_{SM}(\theta) = \frac{1}{2}\,
    \mathbb{E}_{p(x)}\!\left[\|s_\theta(x) - \nabla_x \log p(x)\|_2^2\right],
\]
expand the square and drop the term $\frac{1}{2}\mathbb{E}[\|\nabla_x \log
p(x)\|^2]$, because constant w.r.t. $\theta$:
\[
    J_{SM}(\theta) = \frac{1}{2}\,\mathbb{E}_{p(x)}\!\left[\|s_\theta(x)\|_2^2\right]
    - \mathbb{E}_{p(x)}\!\left[s_\theta(x) \cdot \nabla_x \log p(x)\right] + C.
\]
The cross-term equals:
\[
    \int p(x)\, s_\theta(x) \cdot \nabla_x \log p(x)\,dx
    = \int s_\theta(x) \cdot \nabla_x p(x)\,dx.
\]
Applying integration by parts, with the boundary terms vanishing since $p(x)\to 0$ at
infinity:
\[
    \int s_\theta(x) \cdot \nabla_x p(x)\,dx
    = -\int p(x)\,\nabla_x \cdot s_\theta(x)\,dx
    = -\mathbb{E}_{p(x)}\!\left[\nabla_x \cdot s_\theta(x)\right],
\]
where $\nabla_x \cdot s_\theta(x) = \operatorname{Tr}(\nabla_x s_\theta(x))$
is the trace of the Jacobian. Substituting back:
\begin{equation}
    \label{eq:sm_tractable}
    J_{SM}(\theta) = \mathbb{E}_{p(x)}\!\left[
    \frac{1}{2}\|s_\theta(x)\|_2^2 + \nabla_x \cdot s_\theta(x)
    \right] + C.
\end{equation}
This objective is tractable: it requires only samples from $p(x)$ and network
evaluations, with no ground-truth score. Nevertheless, the cost of computing the Jacobian
trace, which requires $O(d)$ vector-Jacobian products, where $d$ is the data dimension, is prohibitively expensive in high dimensions. Denoising Score Matching
eliminates this cost entirely.

\subsection{Denoising Score Matching}
\label{app:subsec:dsm}

Vincent \cite{vincent2011connection} showed that minimising the Fisher divergence
over the corrupted marginal $p_\sigma(\tilde{x}) = \int p(x)\,q_\sigma(\tilde{x}
|x)\,dx$ is equivalent, up to an additive constant independent of $\theta$, to
minimising the Denoising Score Matching (DSM) objective:
\[
    \mathcal{L}_{\text{DSM}}(\theta)
    = \mathbb{E}_{x \sim p(x),\;\tilde{x} \sim q_\sigma(\tilde{x}|x)}
    \!\left[\left\|s_\theta(\tilde{x},\sigma)
    - \nabla_{\tilde{x}}\log q_\sigma(\tilde{x}|x)\right\|^2\right].
\]

\textit{Proof} Both objectives share the term $\mathbb{E}_{p_\sigma}
[\|s_\theta(\tilde{x})\|^2]$, because $p_\sigma(\tilde{x}) = \int p(x)
q_\sigma(\tilde{x}|x)dx$. The difference lies in the cross-terms. The SM
cross-term is:
\[
    \int p_\sigma(\tilde{x})\,s_\theta(\tilde{x})\cdot
    \nabla_{\tilde{x}}\log p_\sigma(\tilde{x})\,d\tilde{x}
    = \int s_\theta(\tilde{x})\cdot\nabla_{\tilde{x}} p_\sigma(\tilde{x})\,d\tilde{x}.
\]
Since $\nabla_{\tilde{x}} p_\sigma(\tilde{x})
= \int p(x)\,q_\sigma(\tilde{x}|x)\,\nabla_{\tilde{x}}\log q_\sigma(\tilde{x}|x)\,dx$:
\[
    = \iint p(x)\,q_\sigma(\tilde{x}|x)\,s_\theta(\tilde{x})
    \cdot\nabla_{\tilde{x}}\log q_\sigma(\tilde{x}|x)\,dx\,d\tilde{x}
    = \mathbb{E}_{x,\tilde{x}}\!\left[s_\theta(\tilde{x})
    \cdot\nabla_{\tilde{x}}\log q_\sigma(\tilde{x}|x)\right],
\]
which is exactly the DSM cross-term. The residual $\mathbb{E}[\|\nabla_{\tilde{x}}
\log p_\sigma\|^2] - \mathbb{E}[\|\nabla_{\tilde{x}}\log q_\sigma\|^2]$ is
constant in $\theta$, so the two objectives have identical gradients.

For a Gaussian kernel $q_\sigma(\tilde{x}|x)=\mathcal{N}(\tilde{x};\,x,
\sigma^2 I)$, the conditional score is analytic:
\[
    \nabla_{\tilde{x}}\log q_\sigma(\tilde{x}|x)
    = -\frac{\tilde{x} - x}{\sigma^2}
    = -\frac{\epsilon}{\sigma}, \qquad
    \epsilon = \tilde{x} - x \sim \mathcal{N}(0, I).
\]
The network thus learns to predict the normalised noise direction from the
corrupted sample, requiring only paired $(x_0, \tilde{x})$ samples and no
Jacobian computation. This is the foundation for the DSM loss used in both
DDPM (Eq.~\ref{eq:ddpm_noise_loss}) and the CSDI objective
(Eq.~\ref{eq:csdi_objective}).

\subsection{Probability Flow ODE via the Fokker--Planck Equation}
\label{app:subsec:fp_ode}

The Fokker--Planck (FP) equation is the mass-conservation law governing how the
density $p_t(x)$ evolves under an It\^{o} SDE. Here we derive the probability
flow ODE from it.\\

\textbf{Intuition} Treat $p_t(x)$ as a cloud of probability mass. Two forces
reshape it: \textit{drift} $f$ acts like a fluid velocity, transporting mass
along deterministic trajectories and \textit{diffusion} $g\,dW_t$ which injects
randomness, spreading and smoothing the cloud like heat diffusing from a peak.
The FP equation encodes both effects as a single continuity equation.\\

\textbf{The FP equation} For the forward SDE (Eq.~\ref{eq:generalsde}) with
scalar $g(t)$, the marginal density satisfies:
\begin{equation}
    \label{eq:fokker_planck}
    \frac{\partial p_t}{\partial t}
    = -\nabla_x \cdot \bigl(f(x_t,t)\,p_t\bigr)
    + \frac{g(t)^2}{2}\,\Delta_x p_t,
\end{equation}
where $\Delta_x = \sum_i \partial^2/\partial x_i^2$ is the Laplacian. The first
term is advection (transport) by the drift; the second is isotropic diffusion that flattens
peaks and fills troughs.

\textbf{Deriving the probability flow ODE} We seek a deterministic ODE
$dx_t = v(x_t,t)\,dt$ whose marginals match those of the SDE. Such an ODE
satisfies the continuity equation $\partial_t p_t = -\nabla_x\cdot(v\,p_t)$.
Setting this equal to the FP equation requires:
\[
    -\nabla_x\cdot(v\,p_t)
    = -\nabla_x\cdot(f\,p_t) + \frac{g(t)^2}{2}\,\Delta_x p_t.
\]
Using the identity $\Delta_x p_t = \nabla_x\cdot(p_t\,\nabla_x\log p_t)$,
the right-hand side becomes $-\nabla_x\cdot\!\left(f\,p_t -
\frac{g(t)^2}{2}\,p_t\,\nabla_x\log p_t\right)$. Matching the divergence terms:
\begin{equation}
    \label{eq:pf_ode_derived}
    v(x_t,t) = f(x_t,t) - \frac{g(t)^2}{2}\,\nabla_x\log p_t(x_t),
\end{equation}
which is the probability flow ODE stated in Section~\ref{sec:scorediffmodels}.
Replacing the intractable score with $s_\theta(x_t,t)$ yields a Neural ODE
with approximately identical marginals, enabling deterministic sampling and
exact likelihood computation via the instantaneous change-of-variables formula.

[TO EVLAUATE WHETHER TO INCLUDE HEUN'S sampler efficiency]
% \subsection{Heun's Sampler: Efficiency Trade-off}
% \label{app:subsec:heun}

% EDM \cite{karras2022elucidating} integrates the probability flow ODE using
% the second-order Heun method, applied in $\sigma$-space (decreasing from
% $\sigma_{\max}$ to $\sigma_{\min}$). Let $F(x_i,\sigma_i)$ denote the ODE
% slope at $(x_i, \sigma_i)$, computed from $D_\theta$ via the probability flow
% ODE (Chapter~3). A single step from $\sigma_i$ to $\sigma_{i+1}$ (step size
% $h = \sigma_{i+1} - \sigma_i < 0$) proceeds as:
% \begin{enumerate}
%     \item \textit{Euler predictor}:
%     $\hat{x}_{i+1} = x_i + h\,F(x_i,\sigma_i)$.
%     \item \textit{Heun corrector}: evaluate the slope at the predicted point
%     and average:
%     \[
%         x_{i+1} = x_i + \frac{h}{2}\bigl[
%         F(x_i,\sigma_i) + F(\hat{x}_{i+1},\sigma_{i+1})\bigr].
%     \]
% \end{enumerate}
% The corrector costs one additional NFE per step but upgrades the local
% truncation error from $O(h^2)$ (Euler, first order) to $O(h^3)$ (Heun,
% second order), giving a global error of $O(h^2)$.

\section{The Wiener Filter}
\label{app:sec:vp_wiener}

% \subsection{The Variance-Preserving SDE}
% \label{app:subsec:vp_sde}

% The VP SDE from Section~\ref{sec:scorediffmodels}
% uses a linear mean-reverting drift and a state-independent diffusion:
% \[
%     dx_t = -\frac{1}{2}\beta(t)\,x_t\,dt + \sqrt{\beta(t)}\,dW_t,
% \]
% with $\beta(t) > 0$ a monotonically increasing noise schedule. Define
% $B(t) = \int_0^t \beta(s)\,ds$. Since this is a linear SDE, its solution is:
% \[
%     x_t = e^{-\frac{1}{2}B(t)}\,x_0
%     + \int_0^t e^{-\frac{1}{2}(B(t)-B(s))}\sqrt{\beta(s)}\,dW_s.
% \]
% The stochastic integral is Gaussian with variance, by It\^{o}'s isometry:
% \[
%     \int_0^t e^{-(B(t)-B(s))}\beta(s)\,ds
%     = e^{-B(t)}\int_0^t e^{B(s)}\beta(s)\,ds
%     = e^{-B(t)}\bigl[e^{B(t)}-1\bigr]
%     = 1 - e^{-B(t)},
% \]
% where the inner integral evaluates using $\frac{d}{ds}e^{B(s)}=\beta(s)e^{B(s)}$.
% The closed-form marginals are therefore:
% \begin{equation}
%     \label{eq:vp_marginal}
%     x_t \mid x_0 \;\sim\; \mathcal{N}\!\left(
%     e^{-\frac{1}{2}B(t)}\,x_0,\;\bigl(1 - e^{-B(t)}\bigr)\,I\right).
% \end{equation}
% Setting $\bar{\alpha}_t = e^{-B(t)}$ recovers the DDPM \cite{ho2020ddpm} marginal
% $x_t|x_0 \sim \mathcal{N}(\sqrt{\bar\alpha_t}\,x_0,(1-\bar\alpha_t)\,I)$,
% confirming that DDPM is a discretisation of the VP SDE. As $t\to T$ with
% $B(T)$ large, $\bar\alpha_T\to 0$ and $x_T\approx\mathcal{N}(0,I)$, the
% tractable prior. Unlike the VE process (Eq.~\ref{eq:ve}), whose marginal
% variance grows without bound, VP keeps $\operatorname{Var}(x_t)\approx 1$
% when $\operatorname{Var}(x_0)=1$, providing more numerically stable training
% at the cost of lower expressivity.

\subsection{Wiener Filter and the Skip Connection Coefficient}
\label{app:subsec:wiener}

The EDM preconditioning (Eq.~\ref{eq:edm_denoiser}) includes a skip connection
$c_{\text{skip}}(\sigma)\cdot x_t$ that bypasses the neural network. Here $c_{\text{skip}}$ is derived as the optimal linear estimator of $x_0$ given $x_t$,
known as the \textit{Wiener filter coefficient}.\\

Under the VE forward process $x_t = x_0 + \sigma\epsilon$,
$\epsilon\sim\mathcal{N}(0,I)$, we seek the scalar $c$ minimising the mean
squared error of the linear predictor $c\,x_t$:
\[
    \min_c\;\mathbb{E}\!\left[\|c\,x_t - x_0\|^2\right].
\]
Setting the derivative with respect to $c$ to zero (per component):
\[
    \frac{d}{dc}\,\mathbb{E}\!\left[\|c\,x_t - x_0\|^2\right]
    = 2\,\mathbb{E}\!\left[x_t(c\,x_t - x_0)\right] = 0
    \;\implies\;
    c^* = \frac{\mathbb{E}[x_0 \cdot x_t]}{\mathbb{E}[x_t^2]}.
\]
Computing the second moments under $x_t = x_0 + \sigma\epsilon$, with $x_0$
and $\epsilon$ independent and $\sigma_{\text{data}}^2 = \mathbb{E}[x_{0,i}^2]$
the per-component data variance:
\[
    \mathbb{E}[x_0 \cdot x_t]
    = \mathbb{E}[x_0(x_0+\sigma\epsilon)]
    = \mathbb{E}[x_0^2] + \sigma\,\mathbb{E}[x_0\cdot\epsilon]
    = \sigma_{\text{data}}^2,
\]
\[
    \mathbb{E}[x_t^2]
    = \mathbb{E}[(x_0+\sigma\epsilon)^2]
    = \sigma_{\text{data}}^2 + \sigma^2.
\]
Therefore:
\[
    c^* = \frac{\sigma_{\text{data}}^2}{\sigma_{\text{data}}^2 + \sigma^2}
    = c_{\text{skip}}(\sigma).
\]
This is precisely the EDM formula. The interpretation is direct: at low noise
$\sigma \ll \sigma_{\text{data}}$, $c_{\text{skip}}\approx 1$ and the output
is mostly the noisy input (already a good estimate of $x_0$); at high noise
$\sigma \gg \sigma_{\text{data}}$, $c_{\text{skip}}\approx 0$ and the network
$F_\theta$ must reconstruct $x_0$ almost entirely from scratch. The non-linear
residual $c_{\text{out}}(\sigma)\cdot F_\theta(\cdot)$ accounts for what the
Wiener filter cannot: the non-Gaussian, non-linear structure of the data
distribution beyond its second-order statistics.

\section{Bounds on the CRPS}
\label{app:sec:crps_bounds}

Section~\ref{sec:fts_literature} introduced the CRPS as a calibration-adjusted
MAE. Here it is shown that, for $X,X'$ i.i.d.\ draws from $\hat{F}$ and
observation $y$,
\[
    0 \;\le\; \text{CRPS}(\hat{F}, y) \;\le\; \mathbb{E}_{\hat{F}}\!\left[|X-y|\right]
    = \text{MAE}.
\]

\textit{Upper bound.} Since $|X-X'|\ge 0$, its expectation is non-negative, so
\[
    \text{CRPS}(\hat{F}, y)
    = \mathbb{E}\!\left[|X-y|\right] - \tfrac{1}{2}\mathbb{E}\!\left[|X-X'|\right]
    \;\le\; \mathbb{E}\!\left[|X-y|\right] = \text{MAE}.
\]

\textit{Lower bound.} By the triangle inequality, $|X-X'| \le |X-y| +
|X'-y|$. Taking expectations and using that $X$ and $X'$ are identically
distributed, so $\mathbb{E}[|X-y|] = \mathbb{E}[|X'-y|]$:
\[
    \mathbb{E}\!\left[|X-X'|\right]
    \le \mathbb{E}\!\left[|X-y|\right] + \mathbb{E}\!\left[|X'-y|\right]
    = 2\,\mathbb{E}\!\left[|X-y|\right].
\]
Dividing by two and rearranging gives $\mathbb{E}[|X-y|] - \tfrac{1}{2}
\mathbb{E}[|X-X'|] \ge 0$, i.e. $\text{CRPS}(\hat{F}, y) \ge 0$.

Together, the two bounds confirm that CRPS is non-negative and never exceeds
the MAE of the ensemble mean prediction error, with equality in the upper
bound when $\hat{F}$ collapses to a point mass (so $X=X'$ almost surely and
the sharpness term vanishes).